\section{CBT in adult ADHD}\label{sec:adhdcbt}
\subsection{Efficacy}
A randomised trial of 86 medication-treated adults found that CBT outperformed relaxation with educational support, with 53\% versus 23\% responders on clinician global ratings \citep{safren2010cbt}; a smaller earlier trial found 56\% versus 13\% responders for CBT added to medication \citep{safren2005cbt}. A meta-analysis of 32 studies found a between-group effect of $g=.65$ on symptoms, smaller with active controls, and neither medication status nor treatment length moderated outcome \citep{knouse2017meta}. Another found SMD $0.76$ against waitlist but $0.43$ against active controls \citep{young2020cbt}, and a 2026 meta-analysis of 14 RCTs found SMD $-0.45$ on core symptoms and $-0.43$ on executive function \citep{liu2026cbt}. A component network meta-analysis of 113 RCTs found CBT better than placebo only on clinician-reported measures \citep{ostinelli2025adultnma}, and blinded-rater effects of psychological treatments were attenuated to non-significance in an earlier synthesis that did not separate adult from child samples \citep{sonugabarke2013nonpharm}. An umbrella review judged CBT, atomoxetine and methylphenidate to have at least moderate-certainty evidence in adults, with no high-certainty long-term evidence for any intervention \citep{gosling2025umbrella}. The Cochrane review rates the CBT evidence as low quality \citep{lopez2018cochrane}.

\subsection{What the protocols contain, and where CBT sits}
The effective protocols are skills-heavy: time management, organisation and planning, with adaptive thinking as one module \citep{safren2010cbt,solanto2010metacognitive}. One review argued that learning behavioural compensatory strategies ``may be a critical active ingredient'' \citep{knouse2010current}, and the authors of one trial proposed behavioural control as the key mechanism \citep{young2017functional}. We found no dismantling study of cognitive restructuring in adult ADHD, so this is inference. In the largest multi-centre trial ($N=433$), group CBT did not beat structured clinical management, while methylphenidate beat placebo \citep{philipsen2015group}. Guidelines position CBT as an adjunct to medication \citep{kooij2019european,nice2018ng87}. Results on adding medication to CBT are mixed: no incremental benefit in a small double-blind pilot \citep{weiss2012cbt}, but greater early improvement in an unblinded trial ($N=88$), with the gap narrowing over follow-up \citep{cherkasova2020cbtmed}.

\subsection{Anxiety in ADHD}
In a meta-analysis of 13 RCTs ($n=744$), CBT-based interventions reduced anxiety in ADHD modestly (SMD $-0.29$ post-treatment, $-0.42$ at follow-up), with effects more apparent in adults; most trials delivered ADHD-focused rather than anxiety-targeted CBT \citep{wang2026cbtanxiety}. Another meta-analysis found CBT improved anxiety, particularly at follow-up, predicted by core-symptom reduction \citep{lopezpinar2020internalizing}. In children with primary anxiety in group CBT ($N=842$), neither ADHD diagnosis nor subtype predicted response or remission, which cuts against the strong form of our thesis \citep{gould2018cbtanxiety}. The adult version of that question appears untested.
